
\documentclass[11pt,a4paper]{moderncv}
\moderncvtheme[green]{classic}             
\usepackage[T1]{fontenc}
\usepackage[utf8x]{inputenc}
\usepackage[italian]{babel}
\usepackage[scale=0.8]{geometry}
\recomputelengths                             

\fancyfoot{} 
\renewcommand*{\namefont}{\fontsize{25}{28}\mdseries\upshape}

\firstname{Patricio Nicolas}
\familyname{Massaro Rocca}
\mobile{+1-202-754-6052}                    
\email{patomassaro@gmail.com}                

\begin{document}
\maketitle
\vspace{-10mm}

\section{Work experience}

\cventry{2020 - 2024 }{Lead Machine Learning Engineer}{\href{https://www.bairesdev.com/}{Bairesdev}}{}{}{
\begin{itemize}
    \item Designed, trained, improved, and launched machine learning models using AWS, Databricks, Airflow, MLFlow, and SparkML, significantly enhancing client marketing strategies and operational efficiency.
    \item Proposed and implemented data architecture optimizations that directly enhanced machine learning model performance, increasing key performance metrics by 10\%.
    \item Proposed and designed innovative features, including a natural language processing model, and integrated them into the data pipeline, significantly reducing manual workload by 37\%.
    \item Improved model and system performance evaluation using advanced data analysis with Spark SQL and Tableau, supporting strategic decisions among executive leadership.
    \item Collaborated with product and engineering teams to tailor and implement machine learning solutions, significantly boosting client outcomes and aligning with strategic risk management.
\end{itemize} 
}

\cventry{2023-2024}{Data Scientist, Remote Sensing}{\href{https://ororatech.com/}{Ororatech}}{}{}{  
\begin{itemize}
    \item Trained Generative Adversarial Networks (GANs) in Pytorch for remote-sensing applications, achieving a 20\% error reduction in land surface temperature estimation with applications in wildfire management.
    \item Implemented deep convolutional models for georeferencing of thermal data products, enhancing disaster monitoring accuracy.
    \item Published findings at IGARSS 2023 on multi-spectral remote sensing super-resolution.
\end{itemize} 
}

\cventry{2017 - 2020}{Data Engineer}{\href{https://www.tenaris.com/en}{Tenaris}}{}{}{
 \begin{itemize}
     \item Lead the architectural design and workflow development for migrating to a centralized cloud-based PowerBI application, impacting 30,000 users.
     \item Managed data warehousing using Synapse and Data Factory, and designed semantic layers with Azure Analysis Services, reducing costs by 17\%.
 \end{itemize}
}

\cventry{2017 - 2018}{Research Intern}{University of Buenos Aires}{}{}{ 
\begin{itemize}
     \item Authored two publications on the effects of piezoelectric devices in optoacoustic image reconstruction.
 \end{itemize}}

\cventry{2014 - 2017}{Systems Administrator}{Masstec S.R.L}{}{}{
    \begin{itemize}
        \item Administered server, workstation, and networking systems for small businesses.
    \end{itemize}
}
\cventry{2013- 2014}{Network Administrator}{Education Ministry, Buenos Aires}{}{}{
\begin{itemize}
    \item Coordinated the distribution of laptops and integrated ICT into educational practices, enhancing digital literacy for over 250 students.
\end{itemize}
}

\section{Education}
\cventry{2021-2023}{MSc Data Science}{\href{https://www.m-datascience.mathematik-informatik-statistik.uni-muenchen.de/index.html}{Ludwig-Maximilians-Universität}}{Munich}{Germany}{Thesis: Blind Super Resolution Applied to Thermal Remote Sensing Data Products}
\cventry{2020-2021}{MSc Data Mining }{\href{http://datamining.dc.uba.ar/datamining/}{Universidad de Buenos Aires}}{}{Argentina}{}
\cventry{2010-2018}{Electronic Engineering}{\href{https://www.fi.uba.ar/}{Universidad de Buenos Aires}}{}{Argentina}{Thesis: Improvements in Optoacoustic Image Reconstruction Using Deconvolution Filters.} 

\section{Language}
\cvline{Spanish}{Native language.}
\cvline{English}{Bilingual proficiency.}
\cvline{German}{Professional working proficiency.}

\end{document}
